\documentclass[11pt,a4paper]{article}

\usepackage[T1]{fontenc}
\usepackage{lmodern}
\usepackage[margin=21mm,headheight=15pt]{geometry}
\usepackage{amsmath,amssymb,mathtools}
\usepackage{microtype}
\usepackage{xcolor}
\usepackage{enumitem}
\usepackage{fancyhdr}
\usepackage{titlesec}
\usepackage{needspace}

\definecolor{examblue}{HTML}{17365D}
\definecolor{answergreen}{HTML}{245C3A}
\definecolor{rulegray}{HTML}{D6DCE4}

\pagestyle{fancy}
\fancyhf{}
\lhead{\textsc{Machine Learning 1}}
\rhead{\textsc{Mock Exam B -- Answer Key}}
\cfoot{\thepage}
\renewcommand{\headrulewidth}{0.4pt}

\titleformat{\section}
  {\Large\bfseries\color{examblue}}
  {}{0pt}{}
  [\vspace{0.15em}\color{rulegray}\titlerule]
\titleformat{\subsection}
  {\normalsize\bfseries\color{answergreen}}
  {}{0pt}{}

\setlength{\parindent}{0pt}
\setlength{\parskip}{0.4em}
\setlist[enumerate]{leftmargin=1.8em,itemsep=0.8em,topsep=0.35em}
\allowdisplaybreaks

\newcommand{\exercise}[2]{%
  \Needspace{8\baselineskip}%
  \par\addvspace{1.1em}%
  {\Large\bfseries\color{examblue}Exercise #1: #2\par}%
  \nobreak\vspace{0.2em}%
  {\color{rulegray}\hrule height 0.5pt}%
  \vspace{0.55em}%
}
\newcommand{\answer}{\textbf{Answer.}\ }

\begin{document}

\begin{center}
  {\Huge\bfseries\color{examblue} Machine Learning 1}\par
  \vspace{0.35em}
  {\LARGE\bfseries Mock Exam B -- Answer Key}\par
  \vspace{0.7em}
  {\large Worked solutions for all 100 points}
\end{center}

\vspace{0.8em}
\fcolorbox{rulegray}{white}{%
  \begin{minipage}{0.94\textwidth}
    Equivalent algebraic forms are valid. The derivations below use numerator-layout derivatives where a scalar derivative is a row vector.
  \end{minipage}%
}

\exercise{1}{Normalization and a nonsmooth loss}

\subsection*{1. Component derivative}

\answer Let $\rho=\lVert x\rVert_2>0$. Since
\[
  h_i(x)=x_i\rho^{-1},
  \qquad
  \frac{\partial\rho}{\partial x_j}=\frac{x_j}{\rho},
\]
the product rule gives
\[
  \boxed{\frac{\partial h_i}{\partial x_j}
  =\frac{\delta_{ij}}{\rho}-\frac{x_ix_j}{\rho^3}}.
\]

\subsection*{2. Full Jacobian}

\answer With $u=x/\lVert x\rVert_2$,
\[
  \boxed{J_h(x)=\frac{I_d}{\lVert x\rVert_2}
  -\frac{xx^\top}{\lVert x\rVert_2^3}
  =\frac{1}{\lVert x\rVert_2}(I_d-uu^\top)}
  \in\mathbb{R}^{d\times d}.
\]

\subsection*{3. Radial direction}

\answer
\[
  J_h(x)x
  =\frac{x}{\rho}-\frac{x(x^\top x)}{\rho^3}
  =\frac{x}{\rho}-\frac{x\rho^2}{\rho^3}=0.
\]
For $c>0$,
\[
  h(cx)=\frac{cx}{\lVert cx\rVert_2}
  =\frac{cx}{c\lVert x\rVert_2}=h(x).
\]
Thus $h$ is constant along each positive radial ray, so its directional derivative in the radial direction is zero.

\subsection*{4. Absolute-error derivative}

\answer Let $z=a^\top w-y$. For $z\neq 0$,
\[
  \frac{\partial |z|}{\partial w}=\operatorname{sign}(z)a^\top,
  \qquad
  \frac{\partial}{\partial w}\frac{\lambda}{2}w^\top w
  =\lambda w^\top.
\]
Therefore
\[
  \boxed{\frac{\partial f}{\partial w}
  =\operatorname{sign}(a^\top w-y)a^\top+\lambda w^\top}
  \in\mathbb{R}^{1\times p}.
\]

\subsection*{5. Nondifferentiability and subgradient}

\answer At $a^\top w-y=0$, the absolute-value term has different one-sided slopes in direction $a$; the smooth quadratic does not remove this kink. The row-subdifferential is
\[
  \partial f(w)=
  \{sa^\top+\lambda w^\top:s\in[-1,1]\}.
\]
For any chosen $s\in[-1,1]$, one valid column-form update is
\[
  \boxed{w^{+}=w-\eta(sa+\lambda w)}.
\]
For example, $s=0$ gives $w^{+}=(1-\eta\lambda)w$ at the kink.

\exercise{2}{A transformed waiting-time model}

\subsection*{1. Valid density}

\answer The density is nonnegative. With $u=\lambda x^r$ and $du=r\lambda x^{r-1}dx$,
\[
  \int_{-\infty}^{\infty}p(x\mid\lambda)\,dx
  =\int_0^\infty r\lambda x^{r-1}e^{-\lambda x^r}\,dx
  =\int_0^\infty e^{-u}\,du=1.
\]
Thus it is a valid density for every $r,\lambda>0$.

\subsection*{2. Distribution of $Y=X^r$}

\answer For $y>0$, $x=y^{1/r}$ and
\[
  \left|\frac{dx}{dy}\right|
  =\frac{1}{r}y^{1/r-1}.
\]
Hence
\begin{align*}
  p_Y(y\mid\lambda)
  &=r\lambda(y^{1/r})^{r-1}e^{-\lambda y}
    \frac{1}{r}y^{1/r-1}\\
  &=\lambda e^{-\lambda y},\qquad y>0.
\end{align*}
Therefore
\[
  \boxed{Y=X^r\mid\lambda\sim\operatorname{Exponential}(\lambda)},
  \qquad
  \boxed{\mathbb{E}[X^r\mid\lambda]=\frac{1}{\lambda}}.
\]

\subsection*{3. Survival probability}

\answer
\[
  P(X>c\mid\lambda)
  =P(X^r>c^r\mid\lambda)
  =\boxed{e^{-\lambda c^r}}.
\]
For the supplied values,
\[
  P(X>2\mid\lambda=0.25)=e^{-0.25\cdot 2^2}
  =\boxed{e^{-1}}.
\]

\subsection*{4. Maximum-likelihood estimator}

\answer Terms independent of $\lambda$ can be collected into a constant:
\[
  L(\lambda)\propto\lambda^Ne^{-\lambda S},
  \qquad
  \log L(\lambda)=N\log\lambda-\lambda S+\text{constant}.
\]
Thus
\[
  \frac{d}{d\lambda}\log L=\frac{N}{\lambda}-S=0
  \quad\Longrightarrow\quad
  \boxed{\widehat\lambda_{\mathrm{ML}}=\frac{N}{S}}.
\]
The second derivative $-N/\lambda^2$ is negative. Since $S>0$, the log-likelihood also tends to $-\infty$ at both boundaries, so this is the unique global maximum.

\subsection*{5. Posterior distribution}

\answer
\begin{align*}
  p(\lambda\mid x_1,\ldots,x_N)
  &\propto\lambda^Ne^{-S\lambda}
    \lambda^{a-1}e^{-b\lambda}\\
  &=\lambda^{a+N-1}e^{-(b+S)\lambda}.
\end{align*}
Therefore, in shape-rate form,
\[
  \boxed{\lambda\mid x_1,\ldots,x_N
  \sim\operatorname{Gamma}(a+N,b+S)}.
\]

\subsection*{6. MAP estimator}

\answer The derivative of the posterior log-density is
\[
  \frac{a+N-1}{\lambda}-(b+S).
\]
Since $N\geq1$ and $a>0$, the posterior shape is greater than one. Hence the unique interior mode is
\[
  \boxed{\widehat\lambda_{\mathrm{MAP}}
  =\frac{a+N-1}{b+S}}.
\]

\subsection*{7. Posterior mean and asymptotics}

\answer
\[
  \mathbb{E}[\lambda\mid\text{data}]
  =\frac{a+N}{b+S},
  \qquad
  \mathbb{E}[\lambda\mid\text{data}]
  -\widehat\lambda_{\mathrm{MAP}}
  =\boxed{\frac{1}{b+S}}.
\]
If $S/N\to s>0$, the difference is $1/(b+Ns)\to0$ at rate $O(N^{-1})$. Both summaries converge to $1/s$.

\exercise{3}{Bias, variance, and shrinkage}

Write $A=\sigma^2/N>0$ and $D=(\mu-m_0)^2\geq0$.

\subsection*{1. Bias}

\answer
\[
  \mathbb{E}[\widetilde\mu_c]
  =c\mu+(1-c)m_0,
\]
so
\[
  \boxed{\operatorname{Bias}(\widetilde\mu_c)
  =(1-c)(m_0-\mu)}.
\]
The estimator is unbiased exactly when $c=1$ or $m_0=\mu$.

\subsection*{2. Variance}

\answer Independence gives $\operatorname{Var}(\overline X)=\sigma^2/N$. Since $m_0$ is fixed,
\[
  \boxed{\operatorname{Var}(\widetilde\mu_c)
  =\frac{c^2\sigma^2}{N}}.
\]

\subsection*{3. Mean squared error}

\answer
\[
  \boxed{\operatorname{MSE}(c)
  =\frac{c^2\sigma^2}{N}
  +(1-c)^2(m_0-\mu)^2
  =Ac^2+D(1-c)^2}.
\]

\subsection*{4. MSE-minimizing value}

\answer Differentiation gives
\[
  \frac{d}{dc}\operatorname{MSE}(c)
  =2(A+D)c-2D.
\]
Therefore
\[
  \boxed{c^*=\frac{D}{A+D}
  =\frac{(\mu-m_0)^2}
  {(\mu-m_0)^2+\sigma^2/N}}.
\]
This lies in $[0,1)$, and the second derivative $2(A+D)>0$, so it is the unique minimizer.

\subsection*{5. Oracle property and limit}

\answer This is an oracle choice because it depends on the unknown $\mu$. If $\mu\neq m_0$, then $D>0$ and $c^*\to1$ as $N\to\infty$, so the estimator increasingly relies on the precise sample mean. If $\mu=m_0$, then $D=0$ and $c^*=0$ for every $N$; the fixed reference is exactly correct and has zero MSE.

\exercise{4}{Multi-output MAP regression}

Define
\[
  D(W)=\lVert T-\Phi W\rVert_F^2,
  \qquad
  P(W)=\operatorname{tr}((W-M_0)^\top R(W-M_0)).
\]

\subsection*{1. Negative log-posterior}

\answer Up to constants independent of $W$,
\[
  \boxed{-\log p(W\mid T)
  =\frac{\beta}{2}D(W)+\frac{\alpha}{2}P(W)+\text{constant}}.
\]

\subsection*{2. MAP estimator}

\answer Differentiation with respect to $W$ gives the stationarity equation
\[
  \beta\Phi^\top(\Phi W-T)+\alpha R(W-M_0)=0.
\]
With $\lambda=\alpha/\beta$,
\[
  (\Phi^\top\Phi+\lambda R)W
  =\Phi^\top T+\lambda RM_0.
\]
Therefore
\[
  \boxed{W_{\mathrm{MAP}}
  =(\Phi^\top\Phi+\lambda R)^{-1}
  (\Phi^\top T+\lambda RM_0)}.
\]
An equivalent form is
\[
  W_{\mathrm{MAP}}=M_0+
  (\Phi^\top\Phi+\lambda R)^{-1}
  \Phi^\top(T-\Phi M_0).
\]

\subsection*{3. Uniqueness}

\answer For every nonzero $v\in\mathbb{R}^m$,
\[
  v^\top(\Phi^\top\Phi+\lambda R)v
  =\lVert\Phi v\rVert_2^2+\lambda v^\top Rv>0.
\]
Thus the coefficient matrix is positive definite and invertible, regardless of the rank of $\Phi$. The objective is strictly convex, so the MAP estimate is unique.

\subsection*{4. Limits}

\answer If $\Phi$ has full column rank,
\[
  \boxed{\lim_{\lambda\to0}W_{\mathrm{MAP}}
  =(\Phi^\top\Phi)^{-1}\Phi^\top T},
\]
the ordinary least-squares estimate. Using the equivalent form above,
\[
  \boxed{\lim_{\lambda\to\infty}W_{\mathrm{MAP}}=M_0},
\]
because the correction away from $M_0$ is $O(\lambda^{-1})$.

If $\Phi$ is rank deficient, the first displayed inverse does not exist; the small-$\lambda$ limit selects, among least-squares solutions, the one closest to $M_0$ in the $R$-weighted norm.

\subsection*{5. Regularization path and bias--variance tradeoff}

\answer Let $W_i$ minimize $D(W)+\lambda_iP(W)$ for $0<\lambda_1<\lambda_2$, and write $D_i=D(W_i)$ and $P_i=P(W_i)$. Optimality gives
\begin{align*}
  D_1+\lambda_1P_1&\leq D_2+\lambda_1P_2,\\
  D_2+\lambda_2P_2&\leq D_1+\lambda_2P_1.
\end{align*}
Combining these inequalities shows
\[
  \boxed{P_2\leq P_1},
  \qquad
  \boxed{D_2\geq D_1}.
\]
Thus stronger regularization moves the estimate toward $M_0$ and cannot improve the training residual; the inequalities need not be strict.

Under a correctly specified fixed-design model $T=\Phi W_*+E$ with mean-zero isotropic noise, put $A_\lambda=\Phi^\top\Phi+\lambda R$. Then
\[
  \mathbb{E}[W_{\mathrm{MAP}}]-W_*
  =\lambda A_\lambda^{-1}R(M_0-W_*).
\]
Stronger shrinkage generally reduces variance but increases bias toward $M_0$. If $M_0=W_*$, the bias is zero for every $\lambda$, while the variance still decreases.

\exercise{5}{Two sequential Bayesian updates}

Let $\Lambda_N=\Sigma_N^{-1}$.

\subsection*{1. Final precision}

\answer Each Gaussian observation adds a rank-one precision term. Hence
\[
  \boxed{\Lambda_{\mathrm{final}}
  =\Lambda_N+\beta_a\phi_a\phi_a^\top
  +\beta_b\phi_b\phi_b^\top},
\]
and $\Sigma_{\mathrm{final}}=\Lambda_{\mathrm{final}}^{-1}$.

\subsection*{2. Natural parameter and posterior mean}

\answer The linear terms in $w$ give
\[
  \boxed{\eta_{\mathrm{final}}
  =\eta_N+\beta_a\phi_at_a+\beta_b\phi_bt_b}.
\]
Therefore
\[
  \boxed{\mu_{\mathrm{final}}
  =\Sigma_{\mathrm{final}}\eta_{\mathrm{final}}}
\]
or, explicitly,
\[
  \mu_{\mathrm{final}}
  =(\Sigma_N^{-1}+\beta_a\phi_a\phi_a^\top
  +\beta_b\phi_b\phi_b^\top)^{-1}
  (\Sigma_N^{-1}\mu_N+\beta_a\phi_at_a+\beta_b\phi_bt_b).
\]

\subsection*{3. Order independence}

\answer Processing either observation simply adds its precision and natural-parameter terms. Matrix and vector addition are commutative, so both processing orders produce the same final Gaussian posterior. This requires conditional independence:
\[
  p(t_a,t_b\mid w,\phi_a,\phi_b)
  =p(t_a\mid w,\phi_a)p(t_b\mid w,\phi_b).
\]
Equivalently, the two observation noises must be independent conditional on $w$ and the features.

\subsection*{4. Equivalent observation for equal features}

\answer If $\phi_b=\phi_a=\phi$, then the increments are
\[
  \Delta\Lambda=(\beta_a+\beta_b)\phi\phi^\top,
  \qquad
  \Delta\eta=\phi(\beta_at_a+\beta_bt_b).
\]
Define
\[
  \boxed{\beta_*=\beta_a+\beta_b},
  \qquad
  \boxed{t_*=\frac{\beta_at_a+\beta_bt_b}{\beta_a+\beta_b}}.
\]
Then $\Delta\Lambda=\beta_*\phi\phi^\top$ and $\Delta\eta=\beta_*\phi t_*$, exactly the update from one observation $(\phi,t_*,\beta_*)$.

\subsection*{5. Collinear features}

\answer If $\phi_b=c\phi_a$, then
\[
  \beta_b\phi_b\phi_b^\top
  =\beta_bc^2\phi_a\phi_a^\top.
\]
Thus the second observation adds precision only in $\operatorname{span}(\phi_a)$. It provides more information in the same direction, but no new parameter-space direction. If $\phi_a=0$, neither observation gives information about $w$.

\vfill
\begin{center}
  \textbf{End of Mock Exam B Answer Key}
\end{center}

\end{document}
